% !TeX program = xelatex
% !TeX encoding = UTF-8
% SPDX-License-Identifier: MIT
% 修改题名、作者、课程信息，再编辑 sections/example.tex。
\documentclass[coursework]{zzueng}
\usepackage{tikz}
\usetikzlibrary{arrows.meta,positioning,shapes.geometric}

\zzutitle{课程论文排版与实验结果表达示例}
\zzuauthor{你的姓名}
\zzuaffiliation{课程：课程名称\quad 学院：学院名称\quad 学号：你的学号}
\zzuarticleid{} % 需要首页额外显示课程编号时填写，否则留空。

\zzuabstract{课程论文需要将研究问题、分析过程和结论组织为清晰的学术文本。
本示例提供中文题名、作者信息、双栏正文、数学公式、双语图表题和文末英文信息等常用
组成部分，说明如何利用交叉引用维护图表与公式编号，并演示三线表和矢量流程图的写法。
文中的数据由人工构造，用于展示排版结构；参考文献指向公开的排版工具文档。
使用时应将示例内容替换为自己的课程主题、分析过程和实验结果，并按教师要求调整章节、
摘要长度及引用格式。}
\zzukeywords{课程论文；LaTeX；双栏排版；交叉引用；图表表达}

\zzuentitle{An Example of Typesetting a Course Report and Presenting Results}
\zzuenauthor{YOUR Name}
\zzuenaffiliation{Course: Your Course Name; School: Your School; Student ID: Your Student ID}
\zzuenabstract{A course report should organise its question, analysis, and conclusions into
clear academic prose. This example demonstrates a Chinese title block, a two-column body,
mathematical equations, bilingual figure and table captions, and English information at the end
of the report. It also illustrates cross-references, three-line tables, and an editable vector
flowchart. The numerical data are manually constructed for layout demonstration, and the
references point to public typesetting documentation. Students should replace the example
with their own topic, analysis, and results, and follow the requirements of their course.}
\zzuenkeywords{course report; LaTeX; two-column typesetting; cross-reference; visual presentation}

\begin{document}
\maketitle
% SPDX-License-Identifier: MIT
% 本文是为模板新写的教学示例，数据均为人工构造。
课程论文的排版应帮助读者理解内容。写作时可以先交代问题与目标，再说明方法、
展示结果并给出结论。LaTeX 的章节、公式及交叉引用命令可以减少手工调整编号的工作，
相关命令可查阅项目官方文档\cite{latex-doc}；中文字体和标题设置可参考
CTeX 宏集文档\cite{ctex-doc}。

\section{写作与排版方法}
\subsection{组织章节}
将正文写在本文件中，使用标准章节命令组织层次。一级标题用于划分主要部分，
二级标题用于说明具体问题，三级标题用于必要的进一步展开。可以根据课程要求新增
章节，也可以拆分成多个文件并在主文件中引用。

\subsubsection{保持内容与格式一致}
题名、摘要、图表与结论应围绕同一主题。首次出现的术语和变量应解释清楚。
修改章节或图表后重新编译，文中的编号会自动更新。

\subsection{输入数学公式}
以比例指标的写法为例：
\begin{equation}
  p = \frac{m}{n}\times 100\% .
  \label{eq:ratio}
\end{equation}
式中：\(p\)为比例；\(m\)为满足条件的样本数；\(n\)为样本总数，且\(n>0\)。
若人工设定\(m=80\)、\(n=100\)，则\(p=80\%\)。这里的数值仅用于展示
式\eqref{eq:ratio}的排版和引用方式。

正文中的变量使用斜体，单位与说明性下标使用正体，例如\(t_{\mathrm{avg}}\)
表示平均用时，单位为\(\mathrm{ms}\)。向量和矩阵可用\(\bm{x}\)、\(\bm{A}\)表示。
行内简单分式宜写为\(m/n\)，避免影响正文行高。

\section{图表与结果表达}
\subsection{三线表}
表\ref{tab:demo}展示同一指标的比较方式。数据由人工构造，不代表任何算法的实测结果。
在正式作业中，应填写真实结果，并解释统计方式、实验条件和误差来源。

\begin{table}[!htbp]
  \centering\small
  \zzucaption{人工构造的指标比较数据}{Manually constructed comparison data}
  \label{tab:demo}
  \begin{tabular*}{\columnwidth}{@{\extracolsep{\fill}}cccc@{}}
    \toprule
    设置 & 样本数 & 满足条件数 & 比例/\% \\
    \midrule
    A & 100 & 80 & 80.0 \\
    B & 100 & 85 & 85.0 \\
    C & 100 & 90 & 90.0 \\
    \bottomrule
  \end{tabular*}
\end{table}

数据列较多时，可以使用通栏表。表\ref{tab:wide}给出结构示例；表题在表体上方，
中文题后紧接英文题。数值的小数位应与指标精度一致。

\begin{table*}[!t]
  \centering\small
  \zzucaption{不同输入规模下的人工演示数据}{Illustrative data for different input sizes}
  \label{tab:wide}
  \begin{tabular*}{\textwidth}{@{\extracolsep{\fill}}ccccccc@{}}
    \toprule
    \multirow{2}{*}{输入规模} & \multicolumn{3}{c}{用时/\(\mathrm{ms}\)} &
      \multicolumn{3}{c}{比例/\%} \\
    \cmidrule(lr){2-4}\cmidrule(lr){5-7}
      & 设置A & 设置B & 设置C & 设置A & 设置B & 设置C \\
    \midrule
    100 & 12.4 & 10.9 & 9.7 & 80.0 & 85.0 & 90.0 \\
    200 & 24.8 & 21.8 & 19.4 & 80.0 & 85.0 & 90.0 \\
    400 & 49.6 & 43.6 & 38.8 & 80.0 & 85.0 & 90.0 \\
    \bottomrule
  \end{tabular*}
\end{table*}

\subsection{矢量流程图}
图\ref{fig:workflow}由 TikZ 绘制，图形中的文字和连线可以直接修改。
使用自己的图片时，可将文件放入图片目录，并通过插图命令引用。
曲线和流程图宜保留矢量源文件，照片则应保证清晰度。

\begin{figure}[!htbp]
  \centering
  % SPDX-License-Identifier: MIT
\begin{tikzpicture}[
  node distance=6mm,
  font=\fontsize{9bp}{11bp}\selectfont,
  >=Stealth,
  line width=.6pt,
  terminal/.style={draw,rounded corners=3mm,minimum width=22mm,minimum height=7mm},
  process/.style={draw,align=center,text width=45mm,minimum height=8mm},
  decision/.style={draw,diamond,aspect=3,align=center,inner sep=1.5pt,text width=30mm}
]
  \node[terminal] (start) {开始};
  \node[process,below=of start] (requirements) {明确课程要求与主题};
  \node[process,below=of requirements] (write) {编写正文、公式与图表};
  \node[decision,below=of write] (check) {引用与内容完整？};
  \node[process,below=of check] (review) {编译并复核页面};
  \node[terminal,below=of review] (end) {完成};
  \draw[->] (start)--(requirements);
  \draw[->] (requirements)--(write);
  \draw[->] (write)--(check);
  \draw[->] (check)--node[right]{Y}(review);
  \draw[->] (review)--(end);
  \draw[->] (check.east)--++(10mm,0) node[midway,above]{N} |- (write.east);
\end{tikzpicture}

  \zzucaption{课程作业的编写与复核流程}{Writing and reviewing a course report}
  \label{fig:workflow}
\end{figure}

\subsection{讨论结果}
结果讨论应回答研究问题，说明数据支持什么结论，以及结论的适用条件。
本例中的表格仅用于展示结构，因此不据其数字作实际性能判断。正式作业中，
应将图表与相应的分析段落联系起来，并说明异常值或不确定性的处理方式。

\section{结论}
\begin{enumerate}
  \item 章节层次、自动编号和交叉引用可以帮助维护课程论文的结构。
    替换内容后，应检查所有图表和公式引用是否仍指向正确对象。
  \item 图表应配合正文解释结果。示例中的人工数据和流程应根据实际课程任务替换，
    编译后再检查换页、图表尺寸与参考文献。
\end{enumerate}

% SPDX-License-Identifier: MIT
% 两条可核对的公开工具文档；替换为课程作业实际引用。
\begin{thebibliography}{99}
  \bibitem{latex-doc}
  The LaTeX Project. LaTeX documentation[EB/OL]. [2026-10-26].
  \url{https://www.latex-project.org/help/documentation/}.

  \bibitem{ctex-doc}
  CTeX. CTeX: LaTeX classes and packages for Chinese typesetting[EB/OL]. [2026-10-26].
  \url{https://ctan.org/pkg/ctex}.
\end{thebibliography}

% 不要求英文信息的课程，可删除下一行。
\makeenglishabstract
\end{document}
